\section{The two-point skeleton and its moment defect}
\label{ms:geometry-counts}

Every centered law can be separated into its positive and nonpositive
parts. Replacing each part by its conditional mean gives a two-point law;
the remaining conditional fluctuations measure its departure from that
law. We first make this decomposition exact, then identify a moment
defect that vanishes at Esseen's extremal distribution.

\subsection{The sign split and its units}

Fix $P\in\mathcal P$ and put
\[
 p=\Pp(X>0),\quad q=1-p,\quad v=pq,\quad
 m=\E X_+=\E X_->0,\quad B_0=\mathbf1_{\{X>0\}}.
\]
Both $p,q$ are positive; a possible atom at zero is assigned to the
nonpositive class. Define $V=\E[X\mid B_0]$ and $W=X-V$, so that
\[
 V=-m/q\quad(B_0=0),\qquad V=m/p\quad(B_0=1),\qquad
 \theta=\E W^2=1-m^2/v<1.
\]
Write $\theta_-=\E[W^2;B_0=0]$ and $\theta_+=\E[W^2;B_0=1]$.
The distance between the two conditional means is
\[
 \ell=m/v=\sqrt{(1-\theta)/v}>0.
\]
Set $U=W/\ell$, $Y=B_0+U$, and $s=\E U^2$. Then
\begin{subequations}\label{ms:bernoulli-coordinates}
\begin{align}
 X&=\ell(Y-p)=\frac{B_0+U-p}{\sqrt{pq+s}},\qquad
 \E[U\mid B_0]=0,\label{revision:canonical-decomposition}\\
 1&=pq\ell^2+\theta,\qquad s=\theta/\ell^2,\qquad
 v+s=\ell^{-2}.\label{short:acq:normalization}
\end{align}
\end{subequations}
Let $U_0,U_1$ have the two conditional laws, and put
\[
 w_i=\E U_i^2,\quad \mu_i=\E U_i^3,\quad
 s=qw_0+pw_1=\frac{v\theta}{1-\theta},\quad
 \lambda=ns.
\]
The moments $\mu_i$ are finite because $\rho<\infty$. Higher moments
will be used only after bounded support is known: for extremizers this
follows from Proposition~\ref{v9:core:attainment}, and for the local
comparison it is an explicit hypothesis.

If $s=0$, both nonnegative conditional variances vanish, so $U_0=U_1=0$
almost surely. This is exactly the standardized Bernoulli law. If $s>0$,
the parameter $\alpha=pw_1/s$ lies in $[0,1]$ and gives
$w_0=s(1-\alpha)/q$, $w_1=s\alpha/p$.

The signs of $X$ are still those defining the two classes, so expansion
of the actual conditional moments gives
\begin{align}
 \rho/\ell^3&=v(1-2v)+3v(w_0+w_1)-q\mu_0+p\mu_1,\notag\\
 \gamma/\ell^3&=v(q-p)+3v(w_1-w_0)+q\mu_0+p\mu_1,\notag\\
 \beta/\ell^2&=v(q-p)+pw_1-qw_0.
\label{ms:conditional-moments}
\end{align}
For example, on the nonpositive class $|X|^3=\ell^3(p-U_0)^3$;
on the positive class it is $\ell^3(q+U_1)^3$. Conditional centering removes
the linear terms and gives the first identity; the other two follow in
the same way. The three identities share the same pair $(\mu_0,\mu_1)$.

\subsection{A nonnegative moment defect}

The constant $3+\sqrt{10}$ has a simple geometric explanation. The
signed third moment and the first absolute moment satisfy a circle
inequality; the lengths of the support gaps are controlled by the same
first absolute moment. Neither fact requires bounded support.

Define
\[
 \tau(t)=\E[X\mathbf1_{\{X\ge t\}}],\qquad \sigma=\sqrt{10}.
\]
For $t<0$, centering gives $\tau(t)=-\E[X\mathbf1_{\{X<t\}}]\ge0$;
for $t\ge0$ nonnegativity is immediate. Signwise Tonelli integration gives
\begin{equation}
 \int_\R\tau(t)\,dt=1,
 \qquad 2\int_\R|t|\tau(t)\,dt=\rho.
\label{ms:tau-moments}
\end{equation}
On each bounded complementary interval $(a,b)$ of the closed support,
$\tau$ is constant in the interior, with value $\tau_{a,b}$. Put
\begin{equation}
 J=\sum_{(a,b)}\tau_{a,b}(b-a)^2,
 \qquad D=(3+\sigma)\rho-\gamma-3J.
\label{ms:physical-defect-definition}
\end{equation}
There are at most countably many gaps; the next proof shows that the
nonnegative sum is finite. Endpoint atoms do not affect its integrals.

\begin{lemma}[A moment circle and the gap inequality]
\label{ms:physical-defect}
\label{ms:fourier:afcore:lem:physical-defect}
For every $P\in\mathcal P$,
\begin{equation}
 \gamma^2+(2m)^2\le\rho^2,\qquad J\le\rho+2m,
 \qquad \gamma+tJ\le\bigl(t+\sqrt{1+t^2}\bigr)\rho\quad(t>0).
 \label{ms:gap-family}
\end{equation}
The last constant is sharp. Equality holds exactly at the standardized
Bernoulli law with positive mass $p_t=(1+t-\sqrt{1+t^2})/2$.
In particular, $D\ge0$. If zero belongs to the support, then
$D\ge(\sigma-1)\rho$; hence $D<(\sigma-1)\rho$ implies a central
support gap $(-A_-,A_+)$ with $A_-,A_+>0$.
\end{lemma}
\begin{proof}
Write $a_\pm=\E X_\pm^2$ and $b_\pm=\E X_\pm^3$, so
$a_++a_-=1$, $b_++b_-=\rho$ and $b_+-b_-=\gamma$.
Cauchy--Schwarz on each sign class gives
$a_\pm^2\le mb_\pm$. It also gives
$m^2\le pa_+$ and $m^2\le qa_-$, where an atom at zero is included
in $q$. Thus
$m^2(1/a_++1/a_-)\le1$, or $m^2\le a_+a_-$.
Consequently $b_+b_-\ge(a_+a_-)^2/m^2\ge m^2$, proving the
circle inequality since $\rho^2-\gamma^2=4b_+b_-$.

Every gap on one side of zero satisfies
$(b-a)^2\le2\int_a^b|x|\,dx$.
If a central gap $(-A_-,A_+)$ exists, its weight is $m$ and its excess
over this integral is $2mA_-A_+$. Moreover
$0\le\E[(X+A_-)(X-A_+)]=1-A_-A_+$.
Summing the gap bounds by Tonelli proves $J\le\rho+2m$.
If zero is in the support there is no central gap, so $J\le\rho$.
Finally Cauchy--Schwarz in $\mathbb R^2$ gives
$\gamma+2tm\le\sqrt{1+t^2}\rho$, proving the family and the
assertions about $D$.

Equality in the circle forces equality in the two signwise Cauchy
bounds: the positive and negative parts are each concentrated at one
nonzero value. Equality in $m^2\le a_+a_-$ excludes mass at zero.
For such a law, with positive mass $p$ and $q=1-p$, one has
$J=\rho+2m$. Equality in the last Cauchy bound requires
$\gamma=2m/t$, equivalently $t(1-2p)=2p(1-p)$.
Its unique solution in $(0,1)$ is the stated $p_t$.
\end{proof}

\begin{corollary}[The zero-defect law]\label{ms:zero-defect}
For $P\in\mathcal P$, $D=0$ if and only if
\[
 X=(B_0-p_*)/\sqrt{p_*q_*},\qquad
 p_*=(4-\sqrt{10})/2,\quad q_*=1-p_*,
 \quad B_0\sim\operatorname{Bernoulli}(p_*).
\]
\end{corollary}
\begin{proof}
Apply the equality case of Lemma~\ref{ms:physical-defect} with $t=3$.
\end{proof}

\begin{corollary}[The lattice consequence]\label{opus:lattice-gap}
Let $P\in\mathcal P$ be a lattice law with maximal span $h>0$. Then
\[
 J\ge h,\qquad |\gamma|+3h\le(3+\sqrt{10})\rho.
\]
Equality in the second inequality holds exactly for the law of
Corollary~\ref{ms:zero-defect} and its reflection. If successive support
points are all separated by $h$, then $J=h$.
\end{corollary}
\begin{proof}
The support is a closed subset of a translate of $h\mathbb Z$, so its
points have Lebesgue measure zero and every bounded gap has length
$k_gh$ for an integer $k_g\ge1$. The density $\tau$ vanishes outside
the convex hull of the support. Hence, by nonnegative integration,
\[
 \sum_g\tau_gk_gh=\int_\R\tau(t)\,dt=1,
 \qquad J=\sum_g\tau_gk_g^2h^2\ge h.
\]
These identities also hold for unbounded support: the gaps are
countable and their contributions increase to the full integral.
If all $k_g=1$, the inequality is an equality. Under reflection the
zero-bias density becomes $\tau(-t)$ almost everywhere. Thus reflection
preserves $\rho$, $J$ and $h$, and changes the sign of $\gamma$. Apply
Lemma~\ref{ms:physical-defect} to the orientation with nonnegative
third moment. Equality forces $D=0$, and the equality law has
$J=h$, so Corollary~\ref{ms:zero-defect} proves both directions.
\end{proof}
For a lattice law, $J$ is therefore a weighted average of support-gap
lengths, with weights $\tau_g|g|$ summing to one. The same functional
is defined for arbitrary support, which is why it can replace the
lattice span in the defect inequality.

In particular the defect separates two losses:
\[
 D=\underbrace{\sigma\rho-\gamma-6m}_{\text{moment loss}}
       +3\underbrace{(\rho+2m-J)}_{\text{gap loss}},
\]
each nonnegative. The point $(\gamma/\rho,2m/\rho)$ lies in the upper
unit disk. The line $x+3y=\sqrt{10}$ supports that disk at
$(1/\sqrt{10},3/\sqrt{10})$, explaining the coefficient in the moment
loss. For a standardized Bernoulli law with positive mass $p$, direct
calculation gives
\begin{align}
 J&=\frac1{\sqrt{pq}},&
 \rho_B(p)&=\frac{1-2pq}{\sqrt{pq}},\notag\\
 \gamma_B(p)&=\frac{q-p}{\sqrt{pq}},&
 D_B(p)&=\frac{(6+2\sigma)(p-p_*)^2}{\sqrt{pq}}.
 \label{ms:bernoulli-defect}
\end{align}
Indeed its zero-bias density is $\sqrt{pq}$ on the interval between its
two support points. The probability displacement therefore has a
quadratic cost. The equality proof above is what rules out other laws
with zero defect.

For quantitative estimates we need more information than their signs.
Suppose a central gap $(-A_-,A_+)$ has been established,
and write
\[
 T_+=\E X_+^3-mA_+^2,\quad T_-=\E X_-^3-mA_-^2,\qquad
 R_+=\E X_+^2-mA_+,\quad R_-=\E X_-^2-mA_-.
\]
These four quantities are nonnegative. For the outer gaps put
$J_o=J-m(A_-+A_+)^2$. A positive outer gap satisfies
$(b-a)^2\le2\int_a^b(t-A_+)\,dt$, so Tonelli gives
\[
 J_o^+\le2\int_{A_+}^\infty(t-A_+)\tau(t)\,dt
       =\E[X_+(X_+-A_+)^2]=T_+-2A_+R_+\le T_+.
\]
Similarly $J_o^-\le T_--2A_-R_-\le T_-$.
Expanding $D$, using $\sqrt{\sigma+1}\sqrt{\sigma-1}=3$, gives
\begin{align}
 D={}&m\{\sqrt{\sigma+1}A_--\sqrt{\sigma-1}A_+\}^2
       +(\sigma-1)T_++(\sigma+1)T_-\notag\\
    &+3(T_++T_--J_o).\label{ms:defect-decomposition}
\end{align}
Every term is nonnegative. Cauchy--Schwarz also gives
$1/(2\rho)\le m\le1/2$: indeed $1\le\rho\E|X|=2m\rho$ and
$\E|X|\le1$. Since $\sigma-1>2$, the decomposition yields
$T_++T_-\le D/2$. If $D<2$, then
\[
 m(A_-^2+A_+^2)=\rho-T_+-T_-\ge1-D/2,
 \qquad A_-+A_+\ge\sqrt{2-D}.
\]

\subsection{One joint budget for endpoint motion and noise}

Let $\bar A_-=m/q=\ell p$, $\bar A_+=m/p=\ell q$ be the conditional mean
magnitudes. For a central support gap $(-A_-,A_+)$ define
\[
 \Delta_-=\bar A_--A_-,\quad \Delta_+=\bar A_+-A_+,\qquad
 a_0=\Delta_-/\ell,\quad a_1=\Delta_+/\ell,\quad d=a_0+a_1.
\]
These quantities are nonnegative, $d<1$, and
$U_0\le a_0$, $U_1\ge-a_1$. The physical displacement
$\Delta_-+\Delta_+$ is $\ell d$; the displacement $d$ is in count units.
Put $w=1-\theta=pq\ell^2$ and
\[
 \mathcal Q(p)=(\sigma+1)p^2+(\sigma-1)q^2-6pq
       =(6+2\sigma)(p-p_*)^2.
\]

\begin{lemma}[Joint defect budget]\label{ms:joint-defect}
For every $P\in\mathcal P$ whose closed support has a central gap
$(-A_-,A_+)$ with $A_-,A_+>0$,
\begin{align}
 D\ge{}&\ell w\mathcal Q(p)+6\ell w\{d(1-d)+a_0a_1\}\notag\\
 &+\E[\{(\sigma-1)(X+2\bar A_+)+6A_+\}W^2;X>0]\notag\\
 &+\E[\{(\sigma+1)(-X+2\bar A_-)+6A_-\}W^2;X<0].
 \label{ms:joint-defect-budget}
\end{align}
All terms on the right are nonnegative. In order, they charge
$p\ne p_*$, inward motion of the gap endpoints away from their conditional
means, and the spread within each sign class. The last two expectations
retain the actual locations of the conditional fluctuations.
\end{lemma}
\begin{proof}
Write $M_{j,\pm}=\E X_\pm^j$. The outer-gap bounds proved above give
\[
 J_o^+\le M_{3,+}-2A_+M_{2,+}+A_+^2m,\qquad
 J_o^-\le M_{3,-}-2A_-M_{2,-}+A_-^2m.
\]
Adding the central term $m(A_-+A_+)^2$ and substituting in $D$ yields
\begin{align*}
 D\ge{}&(\sigma-1)M_{3,+}+(\sigma+1)M_{3,-}
       +6A_+M_{2,+}+6A_-M_{2,-}\\
       &-6m(A_-^2+A_-A_++A_+^2).
\end{align*}
Conditional centering gives the exact identities
\[
 M_{2,+}=m\bar A_++\E[W^2;X>0],\qquad
 M_{3,+}=m\bar A_+^2+\E[(X+2\bar A_+)W^2;X>0],
\]
and the analogous identities with $-X,\bar A_-$ on the negative class.
After substituting $A_-=\ell(p-a_0)$ and $A_+=\ell(q-a_1)$, the
deterministic terms are
\[
 \ell w\{\mathcal Q(p)+6[d(1-d)+a_0a_1]\};
\]
the remaining terms are exactly the two expectations in
\eqref{ms:joint-defect-budget}. The gap sums are nonnegative, so
monotone convergence justifies summing countably many gaps. Each
expectation is integrable under the finite third-moment assumption;
endpoint atoms do not affect the Lebesgue integrals.
\end{proof}

\begin{corollary}[Simultaneous gap inequality]
\label{ms:simultaneous-gap}
For this central gap,
\begin{align}
 D\ge{}&\frac{(6+2\sigma)(1-\theta)^{3/2}}{\sqrt{pq}}(p-p_*)^2
 +6m\{(A_-+A_+)(\Delta_-+\Delta_+)+\Delta_-\Delta_+\}\notag\\
 &+3(\sigma+1)A_+\theta_++3(\sigma+3)A_-\theta_-.
\label{ms:simultaneous-gap-formula}
\end{align}
\end{corollary}
\begin{proof}
On the positive support $X\ge A_+$ and $\bar A_+\ge A_+$, so the positive
kernel in \eqref{ms:joint-defect-budget} is at least $3(\sigma+1)A_+$.
On the negative support $-X\ge A_-$ and $\bar A_-\ge A_-$, the corresponding
kernel is at least $3(\sigma+3)A_-$. Finally $m=pq\ell$ and
$A_-+A_+=\ell(1-d)$ identify the two deterministic terms with those in
\eqref{ms:simultaneous-gap-formula}.
\end{proof}
